\subsection*{A faithful joint Fox target for the fourth-power cover}

\paragraph{FF.1 (Proved): the marked covering chains and their character components.}
Retain the literal Ferrari subgroup $C=C_2\amalg C_2\amalg C_2$
of $L$ and set
\[
 F_5=C',\qquad D_4=F_5'F_5^4=C''(C')^4,\qquad
 A=F_5/D_4\simeq(\mathbb Z/4)^5,\qquad k=\mathbb F_9.
 \tag{FF.1a}
\]
All subgroups in this statement are closed. The group $F_5$ is free
profinite of rank five by FC.1. Choose an ordered free basis
$(g_1,\ldots,g_5)$ from the cube Schreier construction, and put
$a_j=\overline g_j\in A$. There are $1024$ deck sheets, so
$[C:D_4]=8\cdot1024=8192$, and $D_4$ is free of rank
$1+1024(5-1)=4097$.

The covering of the five-petal bouquet corresponding to $D_4$
has the exact chain sequence
\[
 0\longrightarrow M=H_1(D_4,k)
 \longrightarrow V=k[A]^5
 \stackrel{\partial}{\longrightarrow}I_A
 \longrightarrow0,\qquad
 \partial=(a_1-1,\ldots,a_5-1),
 \tag{FF.1b}
\]
where $I_A$ is the augmentation ideal of $k[A]$. Thus
$\dim_k(M,V,I_A)=(4097,5120,1023)$.
One can first construct this finite graph for the discrete free
group on the chosen basis and then pass to profinite completion.
Its finite-index subgroup is free, and completion preserves its
finite-coefficient abelianization. This proves that (FF.1b)
computes the indicated continuous profinite homology.

Fix the following based path convention. Its column chain
$d(w)\in V$, together with its endpoint $\bar w\in A$, satisfies
\[
 d(uv)=d(u)+\bar u\,d(v),\qquad
 d(g_j)=e_j,\qquad \partial d(w)=\bar w-1.
 \tag{FF.1c}
\]
These equations also give a direct finite construction for every
profinite word. First put $V_0=\mathbb F_3[A]^5$ and use the
homomorphism from the free profinite group to the finite group
$V_0\rtimes A$ which sends $g_j$ to $(e_j,a_j)$.
Its image is precisely
\[
 \{(v,a)\in V_0\rtimes A:\partial v=a-1\}.
 \tag{FF.1d}
\]
Indeed a path to any specified sheet exists, and its chain can
be changed by every cycle in the graph. On $D_4$, the chain
map is the ordinary mod-three homology map. Its kernel is
exactly $D_4'D_4^3$, so (FF.1d) is the actual finite group
$F_5/(D_4'D_4^3)$. Extending the coefficients injectively from
$\mathbb F_3$ to $k$ gives (FF.1c) with values in $V$.
It does not assert that the original group supplies arbitrary
$k$-cycles.

The field $k$ contains all fourth roots of unity, and $1024$
is invertible in it. Consequently the Fourier decomposition is
exact. For a character $\chi:A\to k^\times$, write
\[
 V_\chi=k^5,\qquad
 \partial_\chi=(\chi(a_1)-1,\ldots,\chi(a_5)-1),\qquad
 M_\chi=\ker\partial_\chi.
 \tag{FF.1e}
\]
For each of the $1023$ nontrivial characters, the row
$\partial_\chi$ is nonzero and there is an exact sequence
\[
 0\longrightarrow M_\chi\simeq k^4
 \longrightarrow V_\chi\simeq k^5
 \stackrel{\partial_\chi}{\longrightarrow}k
 \longrightarrow0.
 \tag{FF.1f}
\]
For the trivial character, the boundary map is zero and
$M_1=V_1\simeq k^5$. Thus
$4097=5+1023\cdot4$. The fifth chain coordinate in
(FF.1f) is part of the target and must be retained.

\paragraph{FF.2 (Proved): exact semilinear transport and the inner-action formula.}
An element $\ell\in L$ acts on $F_5$ by its literal conjugation
automorphism $\phi$. Since $D_4$ is characteristic in $F_5$,
there is an induced automorphism $\bar\phi$ of $A$. Define
the marked Fox Jacobian by its columns:
\[
 J_\phi=(d(\phi(g_1))\ \cdots\ d(\phi(g_5)))
       \in\operatorname{Mat}_5(k[A]).
 \tag{FF.2a}
\]
The based chain substitution is the semilinear map
\[
 T_\phi(v)=J_\phi\,\bar\phi(v),\qquad
 T_\phi(av)=\bar\phi(a)T_\phi(v).
 \tag{FF.2b}
\]
Equation (FF.1c), first on the basis and then on products and
inverses, proves
$d(\phi(w))=J_\phi\bar\phi(d(w))$. Consequently
\[
 \partial J_\phi=\bar\phi(\partial),\qquad
 J_{\phi\circ\psi}=J_\phi\bar\phi(J_\psi),\qquad
 T_{\phi\circ\psi}=T_\phi T_\psi.
 \tag{FF.2c}
\]
The inverse automorphism supplies the inverse chain map.
These are genuine actions of $L$ on the finite chain spaces,
obtained from its source conjugation action. All defining
relators of $L$ therefore hold; there is no assumption that
a gentle element acts on a larger five-point group.
Continuity follows directly from the finite construction
(FF.1c), since each Jacobian column and induced deck action
depends continuously on the images of the finitely many
chosen generators.

For clarity, the exact character transport rule is as follows.
Let
\[
 e_\chi=\frac1{|A|}\sum_{a\in A}\chi(a)^{-1}a.
\]
Then $\bar\phi(e_\chi)=e_{\chi'}$, where
\[
 \boxed{\quad \chi'=\chi\circ\bar\phi^{-1},\qquad
 T_\phi:V_\chi\longrightarrow V_{\chi'},\quad
 v\longmapsto J_\phi(\chi')v.\quad}
 \tag{FF.2d}
\]
The Jacobian is evaluated at the \emph{target} character.
In particular boundary compatibility is the exact row identity
\[
                \partial_{\chi'}J_\phi(\chi')=\partial_\chi.
 \tag{FF.2e}
\]
The literal positive and inverse Nielsen words and Fox Laurent
coefficients are stored in
\path{agent_lifting/records/f5_nielsen_fox.json}.
The independent certificate
\path{agent_lifting/code/check_f5_fox_characters.py} checks (FF.2e) for
all $1024$ characters and all six marked generator maps,
as well as both orders of their inverse products. Its saved
output is \path{agent_lifting/records/check_f5_fox_characters.out}.

No commutativity of the pure-input character permutations is
used here. In fact those permutations need not commute modulo
four. An individual value $f(X_i,X_j)$ must be multiplied with
its actual block permutation, unless its action on $A$ has
separately been proved trivial. Restriction to one character
without that proof would not define the required word
evaluation. The direct sum over all nontrivial characters is
$L$-stable and retains every transition in (FF.2d).

For later faithfulness, consider an actual element $c\in F_5$.
Conjugation by $c$ acts trivially on the abelian deck group $A$.
Using (FF.1c) gives, for each basis element,
\[
 d(cg_jc^{-1})
 =d(c)+\bar c e_j-\bar g_jd(c)
 =\bar c e_j+(1-a_j)d(c).
\]
Therefore its full chain action is
\[
 \boxed{\quad T_{\operatorname{Inn}(c)}
          =\bar c I-d(c)\partial,\qquad
 T_{c,\chi}=\chi(\bar c)I-d(c)_\chi\partial_\chi.\quad}
 \tag{FF.2f}
\]
In particular its action on $M_\chi$ is the scalar
$\chi(\bar c)$. If $c\in D_4$, this becomes
$I-d(c)_\chi\partial_\chi$. It is the identity on
$M_\chi$ for every such $c$, whereas on the full $V_\chi$
it is the identity if and only if $d(c)_\chi=0$, for
$\chi\ne1$. The latter equivalence follows because
$\partial_\chi$ is a nonzero row over a field.

\paragraph{FF.3 (Proved): the joint target detects the actual fourth-power affine quotient.}
Put
\[
 E_0=C/(F_5'F_5^3)=C/(C''(C')^3),\qquad
 W=\bigoplus_{\chi\ne1}V_\chi.
\]
The old affine target is the actual conjugation action on
$E_0$, not merely its rank-five homology action. Together
with (FF.2), it gives a continuous finite representation
\[
 \rho_{\rm joint}:L\longrightarrow
       \operatorname{Aut}(E_0)\times\operatorname{GL}_k(W).
 \tag{FF.3a}
\]
Its kernel on the full Ferrari subgroup $C$ is exactly
\[
 \boxed{\quad C\cap\ker\rho_{\rm joint}=D_4'D_4^3.\quad}
 \tag{FF.3b}
\]
Thus the image of $C$ in this joint target is a faithful copy
of the actual group $C/(D_4'D_4^3)$.
All the quotients here are legitimate: $F_5$ is characteristic
in the normal subgroup $C$, and $D_4'D_4^3$ is characteristic
in the characteristic subgroup $D_4$ of $F_5$, hence normal
in $L$.

Here is a complete proof. First restrict to $F_5$. Suppose
$c\in F_5$ is killed by the
joint target. For each nontrivial character, its restriction
to the four-dimensional space $M_\chi$ is multiplication
by $\chi(\bar c)$, by (FF.2f). The identity action therefore
implies $\chi(\bar c)=1$ for every nontrivial $\chi$.
These characters separate $A$, so $c\in D_4$.
The full five-dimensional blocks now imply
\[
                         d(c)_\chi=0\qquad(\chi\ne1).
 \tag{FF.3c}
\]

It remains to detect the trivial character. Write
$M_0=F_5/(F_5'F_5^3)\simeq\mathbb F_3^5$ and
$A_0=C/F_5\simeq(C_2)^3$. The cube calculation FC.2
decomposes $M_0\otimes k$ into the nontrivial characters
$110,101,011,111,111$ of $A_0$. In particular
\[
                              M_0^{A_0}=0.
 \tag{FF.3d}
\]
If $c\in F_5$ acts trivially by conjugation on $E_0$, its
image in the normal abelian subgroup $M_0$ must be fixed
by $A_0$, hence is zero. Thus the old affine component
detects exactly the ordinary mod-three abelianization
of $c$ in $F_5$. That abelianization is
\[
             d(c)_1=(\text{exponent sums of }c\bmod3)
                         \in k^5.
 \tag{FF.3e}
\]
It follows that this last Fourier component also vanishes.
Fourier inversion in $k[A]$ now gives $d(c)=0$.
Since $c\in D_4$, (FF.1c) identifies its kernel as
$D_4'D_4^3$. Conversely, an element of $D_4'D_4^3$
has endpoint one and every chain component zero, and
also lies in $F_5'F_5^3$; it is therefore killed by both
factors in (FF.3a). This proves the asserted kernel equality
on $F_5$.

To extend it to $C$, the same old affine group is sufficient.
The four characters $110,101,011,111$ span $A_0^*$, so the
$A_0$-action on $M_0$ is faithful. A central element of $E_0$
first has trivial image in $A_0$, because it acts trivially
on $M_0$; it then belongs to $M_0^{A_0}=0$. Thus
$Z(E_0)=1$. If $c\in C$ is killed by the old conjugation
target, its image in $E_0$ is central and hence one. Therefore
$c\in F_5'F_5^3\subset F_5$, where the kernel equality is
already proved. This establishes (FF.3b) on all of $C$.

The trivial component can equivalently be described by
transfer. Let $p_*$ be the map on homology induced by
the $1024$-sheeted covering, and let $\operatorname{tr}$
be transfer. Then
\[
 p_*\operatorname{tr}=1024\,\mathrm{id},\qquad
 \operatorname{tr}p_*=\sum_{a\in A}a.
 \tag{FF.3f}
\]
Because $1024$ is a unit in characteristic three,
$p_*:H_1(D_4,k)^A\to H_1(F_5,k)$ is an isomorphism.
This is precisely the rank-five Fourier component in
(FF.3e). The full chain action at the trivial character
alone cannot detect it: (FF.2f) there is always the
identity. The old affine factor in (FF.3a) supplies
exactly this missing information.

Finally let $f$ be a genuine gentle multiplier-one word,
and use the literal five pentagon inputs $X_i$. Its
pentagon defect $\mathcal P(f)$ belongs to $C'$.
For completeness, this weaker conclusion of RP.8
already follows in $L/C'$: every pure $X_i$ centralizes
$C/C'$, the images of $X_4,X_5$ belong to that abelian
kernel, and $X_3X_1^{-1}$ also belongs to it.
The three nonharmonic commutator-word evaluations are
therefore one, while the other two cancel by the
two-cycle relation. Consequently (FF.3b) gives the
exact equivalence
\[
 \boxed{\quad
 \rho_{\rm joint}(\mathcal P(f))=1
 \quad\Longleftrightarrow\quad
 \mathcal P(f)\in D_4'D_4^3
 \quad\Longleftrightarrow\quad
 \mathcal P(f)=1\text{ in }L/(D_4'D_4^3).\quad}
 \tag{FF.3g}
\]
The stronger previous result $\mathcal P(f)\in
C''(C')^2$ is compatible with this statement but is
not needed for its proof. The equivalence concerns
the product with its actual semilinear transport;
it does not assert that its individual factors are
blockwise or trivial. It establishes a faithful
test in an actual finite relative quotient, without
asserting that the test holds for every gentle
element or proves the full GT sign comparison.

\paragraph{FF.4 (Proved): a conditional transfer to a genuine relative source automorphism.}
Let
\[
 N=D_4'D_4^3,\qquad \Gamma=L/N,\qquad K=C/N,
 \qquad \pi:\Gamma\longrightarrow\widehat B.
 \tag{FF.4a}
\]
For a genuine gentle multiplier-one pair $(1,f)$, put
\[
 g=f(X_1,X_2)=f(a^2,b^2),\qquad
 h=f(X_4,X_3)=f((ab)^3,c^2).
\]
Assume exactly the following two hypotheses in the full joint
chain-and-affine target of FF.3:
\[
             \rho_{\rm joint}(g)=1,
             \qquad \rho_{\rm joint}(h)=1.
 \tag{FF.4b}
\]
Then the literal Drinfeld generator assignment
\[
        a\longmapsto a,\qquad
        b\longmapsto g^{-1}bg,\qquad
        c\longmapsto h^{-1}ch
 \tag{FF.4c}
\]
defines a genuine continuous automorphism of $\Gamma$, fixes
$K$ pointwise, and induces the given gentle automorphism on
$\widehat B$. This is a conditional assertion about an actual
relative quotient; no claim is made that (FF.4b) holds for a
particular proposed seed.

To prove it, let $q:L\to\Gamma$ be the quotient map. By FF.3,
$N\subset\ker\rho_{\rm joint}$, so the joint representation
factors as $\rho_{\rm joint}=\theta q$, where
$\theta:\Gamma\to\operatorname{im}\rho_{\rm joint}$ is
faithful on $K$. Define (FF.4c) first as a map from the free
profinite group on the three braid generators into $\Gamma$.
By (FF.4b), its composite with $\theta$ agrees on every
generator with the original quotient map.

Its fold is also prescribed. In $\widehat B$, the first
argument of $h$ is one, and $f\in\widehat F_2'$, so $h$
has trivial fold. The three folded images are therefore
$a,\,\bar g^{-1}b\bar g,\,a$. These are the images of
$a,b,a$ under the genuine three-strand gentle automorphism,
whose construction uses only the gentle two- and three-cycle
relations and its free-group automorphism, as in FP.26 and
GD.2. In particular the folded map is that automorphism
composed with the original fold.

Consider each defining relator of $L$:
\[
             aba=bab,\qquad bcb=cbc,\qquad
             [a,c]=1,\qquad (abc)^4=1.
\]
Its new image has trivial fold and therefore belongs to $K$.
Its new $\theta$-value is also trivial, because the new and
old generator values under $\theta$ agree. Faithfulness of
$\theta$ on $K$ makes the new relator image one. This proves
all four relators, including the central full-twist relation,
and hence constructs a continuous homomorphism
$\widetilde\Phi:L\to\Gamma$.

For every $w\in C$, both $\widetilde\Phi(w)$ and $q(w)$
belong to $K$, and their $\theta$-values agree. The same
faithfulness gives
\[
                         \widetilde\Phi(w)=q(w)
                                  \qquad(w\in C).
 \tag{FF.4d}
\]
In particular $\widetilde\Phi$ kills $N$. Kernel stability
has thus been proved before descending to a homomorphism
$\Phi_\Gamma:\Gamma\to\Gamma$; no lift of an automorphism
of a deck group to its free covering group has been assumed.
The descended map is the identity on $K$ and the given
automorphism on $\widehat B$. It is injective because an
element of its kernel first has trivial fold and then lies
in the kernel of the identity on $K$. It is surjective by
lifting a desired folded preimage and correcting its residual
error by an element of $K$. Its inverse is continuous because
$\Gamma$ is compact Hausdorff. This proves the assertion.

The two hypotheses (FF.4b) concern full word evaluations in
the semilinear target. Triviality only on the four-dimensional
cycle blocks, a solution of the pentagon equation alone, or
nilpotence in a discarded linear quotient does not supply
these hypotheses. In particular this proposition does not
assert a full GT action, a new character comparison, or an
action on an additional homology cover whose kernel stability
has not been established.

\paragraph{FF.5 (Proved): the primitive-character translation left after the previous affine stage.}
Put $D_2=F_5'F_5^2=C''(C')^2$. For every actual word
$w\in D_2'$, the full joint target has the following exact
reduction. First $w\in F_5'\subset D_4$, so its deck image
in $A$ and its old affine image in $E_0$ are trivial. Moreover
\[
                    d(w)_\chi=0\qquad(\chi^2=1).
 \tag{FF.5a}
\]
There are exactly $32$ such characters, including the trivial
one. To prove (FF.5a), observe that each of them kills $D_2$.
The crossed identity (FF.1c), specialized at such a character,
therefore makes $d_\chi:D_2\to k^5$ an ordinary additive
homomorphism. It kills the closed derived subgroup $D_2'$.

There remain exactly $4^5-2^5=992$ primitive order-four
characters. On each of them write $m_\chi=d(w)_\chi$. The
actual full-chain matrix is
\[
 T_{w,\chi}=I-m_\chi\partial_\chi,\qquad
 \partial_\chi m_\chi=0,\qquad
 \operatorname{rank}(T_{w,\chi}-I)\leq1,\qquad
 (T_{w,\chi}-I)^2=0.
 \tag{FF.5b}
\]
These assertions follow directly from (FF.2f) and
$\partial d(w)=0$. The matrix determines the translation:
for any $j$ with $\chi(a_j)\ne1$,
\[
             m_\chi=
 \frac{e_j-T_{w,\chi}e_j}{\chi(a_j)-1}.
 \tag{FF.5c}
\]
Thus $w=1$ in $C/(D_4'D_4^3)$ if and only if all these
$992$ vectors vanish. No further affine factor is needed
for this restricted class of words. Since the original
chains are defined over $\mathbb F_3$, conjugate characters
satisfy $m_{\chi^3}=m_\chi^3$ coefficientwise. The $496$
conjugate pairs each supply four coordinates over
$\mathbb F_9$, hence at most $3968$ coordinates over
$\mathbb F_3$ in total.

There is also an exact structural meaning to these counts.
In $H_1(D_4,k)$, the characters of order at most two have
combined dimension $5+31\cdot4=129$. They form the
invariant subspace for
\[
                    B=D_2/D_4=2A\simeq(C_2)^5.
\]
Transfer for the covering of index $32$ identifies this
subspace with $H_1(D_2,k)$. Its complementary primitive
part has dimension $4097-129=3968$.
In fact the image of $D_2'$ in $D_4/(D_4'D_4^3)$ is
\emph{exactly} this primitive part over $\mathbb F_3$.
For an elementary verification, consider the finite extension
\[
 1\longrightarrow M_3=D_4/(D_4'D_4^3)
 \longrightarrow D_2/(D_4'D_4^3)
 \longrightarrow B\longrightarrow1.
\]
It splits because $|B|=32$ is invertible on the elementary
abelian three-group $M_3$. Explicitly, for a section with
additive factor set $\alpha(g,h)$, put
$b(g)=|B|^{-1}\sum_h\alpha(g,h)$. Summing the cocycle
identity over its final argument gives
$\alpha(g,h)=b(g)+g b(h)-b(gh)$, and replacing the section
by $(-b(g))$ times that section kills its factor set.
The split extension is therefore $M_3\rtimes B$. As $B$
is abelian, its derived subgroup is precisely $[M_3,B]$,
the sum of the nontrivial $B$-character components. Those
are exactly the primitive $A$-character components above.
The image of $D_2'$ is that derived subgroup, proving the
claim. This surjectivity concerns arbitrary words in
$D_2'$; it does not assert that every such vector arises
as the pentagon defect of a genuine gentle element.

In particular, when the actual witness retained in RA.5
has $\mathcal P(f)\in D_2'$, all lower-order components
are already zero, and (FF.5b)--(FF.5c) isolate its remaining
fourth-power affine pentagon test. They do not apply to an
individual factor in a proposed seed's normal closure merely
because that factor was used to construct a pentagon.
The separately computed product of two seed-conjugate
involutions with order nine on full chains, and a size-five
unipotent Jordan block, is a factor-level normal-closure
obstruction. It is not the shear (FF.5b) of an actual
pentagon known to lie in $D_2'$. The latter must have the
stated rank-one, square-zero form on every primitive block.

\paragraph{FF.6 (Proved): uniform odd-coefficient faithfulness and the exact older-cover factor.}
Let $\ell$ be any odd prime and $n\geq1$. Choose the unramified
coefficient ring $\mathcal O_\ell$ containing all fourth roots
of unity: one can take $\mathbb Z_\ell$ when $\ell\equiv1$
modulo four, and $\mathbb Z_\ell[i]$ when $\ell\equiv3$
modulo four. Put $R=\mathcal O_\ell/\ell^n$ and
\[
 N_{\ell,n}=D_4'D_4^{\ell^n},\qquad
 E_{0,\ell,n}=C/(F_5'F_5^{\ell^n}).
\]
The full Fox representation over $R$, together with the old
affine conjugation action on $E_{0,\ell,n}$, satisfies
\[
 C\cap\ker\rho_{\ell,n}=N_{\ell,n},\qquad
 Z\bigl(C/N_{\ell,n}\bigr)=1.
 \tag{FF.6a}
\]
This includes prime powers, not just residue fields.

For faithfulness, repeat the path construction first over
$\mathbb Z/\ell^n$ and then extend scalars to $R$. The
extension is injective and faithfully flat. Fourier
idempotents exist because $1024$ is a unit. For a nontrivial
character, at least one entry of $\partial_\chi$ is a unit:
distinct fourth roots remain distinct modulo every odd
prime. Thus $\ker\partial_\chi$ is free of rank four,
and $m\partial_\chi=0$ implies $m=0$ over $R$ as well.
The old cube representation has the characters
$110,101,011,111,111$ over $\mathbb Z[1/2]$; this follows
directly by decomposing its signed edge modules with the
eight character idempotents. It has no invariants and its
$(C_2)^3$-action is faithful at every odd modulus. Hence
$E_{0,\ell,n}$ is centreless by the argument in FF.3.
Every step of that proof consequently applies and proves
the kernel equality in (FF.6a).

For the centre assertion, a central element first maps to
one in $E_{0,\ell,n}$, so it has a representative in $F_5$.
Its action on the nontrivial Fourier cycle blocks is the
scalar $\chi(\bar c)$, and centrality forces all these
scalars to be one. Its deck coordinate in $A$ is therefore
one, so the element belongs to
$M_{\ell,n}=D_4/N_{\ell,n}$. An invariant of this module
under $F_5/D_4=A$ belongs to the trivial Fourier component.
Transfer identifies that component with
$H_1(F_5,\mathbb Z/\ell^n)$, equivariantly for $C/F_5$.
The old cube representation has no $C/F_5$-invariants,
so $M_{\ell,n}^{C/D_4}=0$. This proves centrelessness.
Scalar extension was faithful, so the conclusion holds in
the actual group over $\mathbb Z/\ell^n$.

There is an exact comparison for the $32$ characters of
order at most two. Write $A_2=F_5/D_2\simeq(C_2)^5$.
The group-ring map $R[A]\to R[A_2]$ sends the full
marked Fox Jacobians for the fourth-power cover to the
full Jacobians for the second-power cover. Under Fourier
transform this is precisely the direct sum of the
order-at-most-two blocks, including their augmentation
coordinates. Its based path map has kernel
\[
                         D_2'D_2^{\ell^n}.
 \tag{FF.6b}
\]
Consequently, if a literal source automorphism acts as the
identity on $C/(D_2'D_2^{\ell^n})$, it acts as the identity
on every one of these full chain blocks at that modulus.
Indeed for each free generator $g_j$, its image differs
from $g_j$ by an element of $D_2'D_2^{\ell^n}$,
which has zero path chain modulo (FF.6b); its deck action
on $A_2$ is also the identity. Both its Jacobian and its
semilinear coefficient action are therefore the identity.
In particular the low-order full-chain representation of
$L$ factors through its actual conjugation action on
$C/(D_2'D_2^{\ell^n})$, and hence through the faithful
older affine action used in RA.5. The old affine factor
$E_{0,\ell,n}$ also factors
through this same action, since
$D_2'D_2^{\ell^n}\subset F_5'F_5^{\ell^n}$.
This is an exact comparison of based paths and source
kernels, not merely a transfer isomorphism of cycles.
It is therefore applicable to the normal nilpotent source
conditions retained in RA.5, as well as to the pointwise
kernel-fixing conclusion of RA.6.

\paragraph{FF.7 (Proved using R2 and the stated exact certificates): one actual witness passes the full fourth-power abelian family.}
There is one genuine gentle multiplier-one element $k$ such
that
\[
 \boxed{\quad \Delta(k)=1,\qquad D(k)=1,\qquad
       \mathcal P(f_k)=1\ \text{in}\ L/D_4'.\quad}
 \tag{FF.7a}
\]
It simultaneously retains all the earlier conclusions of
LI.8, LI.9, FR.6--FR.7, and RA.5--RA.6. The new quotient
retains $D_4/D_4'\simeq\widehat{\mathbb Z}^{4097}$ in its
entirety. It does not replace this abelianization by the
nonabelian pro-three completion of $D_4$.

We give the source and compactness argument, including the
exact finite inputs. Use the same actual three-supported
cyclic seed $S$ as RA.5, and let
\[
         \mathcal N_0=\overline{\langle\!\langle S
                              \rangle\!\rangle}_{\widehat B}.
\]
It lies in the original orbifold cores, in $\widehat F_2''$,
and in every previously retained normal source condition
that contains $S$. Its signed image is
$q_{12}(\mathcal N_0)=\langle[5I]\rangle\simeq C_2$.
The global fourth-power certificate of FP.56, with the
independent simple-group proof in DF.2, includes all five
pentagon pairs and all six peripheral markings, hence all
thirty source maps. It checks that the images of $S$ fix
character indices, act identically on boundary and
augmentation, and have nontrivial primitive proper support
only in $40$ character orbits for one harmonic pair: $32$
of size eight and eight of size sixteen. The other harmonic
pair is its explicitly checked $abc$-conjugate; both pairs
and all six markings are included in the simultaneous target.
Every active proper normal image
is a certified centreless nonabelian simple group of order
$6048$. The reference group is proved simple by enumeration
of its fourteen conjugacy classes and their normal closures;
all forty images are explicitly conjugated to that reference.
These precise inputs are saved in
\path{agent_five/d4_global_correction_inputs.json},
\path{agent_five/d4_proper_simple_conjugacies.json}, and
\path{agent_lifting/records/d4_proper_stabilizer.json}.
Perfection or group order alone is not substituted for
simplicity in this argument.

Let $E$ be the simultaneous image of $\mathcal N_0$ on
\emph{all} primitive full chain blocks, under these thirty
source maps. Since $\mathcal N_0$ is $B$-normal, the six
markings do not enlarge its possible source values; they
are nevertheless retained explicitly to make all subsequent
kernel conditions $B$-normal. The simultaneous proper image
is a subdirect product of centreless nonabelian simple groups,
so it is a product of diagonal simple factors and is perfect.
To justify this last implication, apply Goursat inductively:
a new simple factor is either independent or identified with
one existing simple factor. A quotient of a product of
nonabelian simple groups onto a simple group factors through
one factor, since the different normal images commute.

The kernel from $E$ to this simultaneous proper image is a
three-group. Every element in the kernel fixes the character
indices and acts identically on all three graded terms of
the boundary/cycle/chain filtration. In an adapted basis it
is a block upper-unitriangular matrix over $\mathbb F_9$;
a product of these finite unipotent groups is a three-group.
This also includes primitive blocks where the proper image
is trivial. It follows that $E^{\rm ab}$ is a three-group,
and hence that $E$ has no nontrivial quotient $C_2$.

Let $\tau:\mathcal N_0\to E$ denote this actual simultaneous
source homomorphism. Its kernel still maps onto
$\langle[5I]\rangle$ under $q_{12}$. Otherwise the nontrivial
quadratic quotient of $\mathcal N_0$ would factor through
$E$, which is impossible. Choose an actual element
\[
       A\in\ker\tau\subset\mathcal N_0,
                    \qquad q_{12}(A)=[5I].
 \tag{FF.7b}
\]
Thus $A$ kills the \emph{whole} primitive full-chain image
modulo three, including its shears, and inherits all the
older normal source conditions. This is an existence
argument in the genuine compact source group; no independent
assignment of pentagon factors is being lifted.

At the prime three, take the common inverse image of the
identity in the primitive chain representation modulo three,
including its six-marking $B$-normal core. It contains $A$.
Its image at every higher three-adic level is contained in
one normal pro-three congruence kernel: successive matrix
congruence kernels are additive groups in characteristic
three. This is a statement about a source subgroup and its
normal image, not the insufficient assertion that the single
element $A$ has three-power order. The low-order blocks and
the old affine factor are quotients of the old $C/D_2'$ action
by FF.6. Their common normal images are the pronilpotent
ones already established in RA.5. The full new image embeds
in the product of that old pronilpotent image and the new
pro-three congruence image. A closed subgroup of a product
of pronilpotent groups is pronilpotent. At each finite
level, therefore, the old and new requirements give one
common normal nilpotent pair-image subgroup whose source
core contains $A$.

For every odd prime $\ell\ne3$, the exact Gaussian-integer
Fox calculation gives
\[
       (X_i^{m_i}-I)^2=0,\qquad
       m_i=4\ (i=1,2,3),\qquad m_i=8\ (i=4,5),
 \tag{FF.7c}
\]
including the actual character permutations. Every third
peripheral $(X_iX_j)^{-1}$ has the same property with a
power of two at most eight. These identities are checked
on all $1024$ full five-dimensional blocks over
$\mathbb Z[i]$, together with all $L$ relators, inverses,
and boundary identities. The exact input is
\path{agent_lifting/records/d4_integral_peripherals.json}
and its constructing certificate
\path{agent_lifting/code/check_d4_integral_peripherals.py}.
At modulus $\ell^n$, (FF.7c) bounds each peripheral order
by $m_i\ell^n$, since $(I+N)^{\ell^n}=I+\ell^nN$ when
$N^2=0$. The three-supported exponent $e_3$ therefore
kills every peripheral. All the values of $S$, and hence
of its $B$-normal closure $\mathcal N_0$, are trivial in
these full chain targets and all their marked cores.
In particular $A$ satisfies every such condition.

At the prime two there is no Fourier division by a power
of two. The group
$C'/(D_4'D_4^{2^n})$ is a finite two-group, so
\[
                    O^2(C')\subset D_4'D_4^{2^n}.
 \tag{FF.7d}
\]
The retained FR.6--FR.7 source conditions in $L/O^2(C')$
therefore imply the required equations and factorwise
centralizations at every two-primary level.

For any finite collection of these prime-power conditions,
use $A$ in the same R2 frame lift and normalization as LI.8
and RA.5. Every prescribed $B$-normal source core contains
$A$, its signed value is still $[5I]$, and the required
normal pair images are nilpotent. At this finite stage,
R2 supplies the prescribed finite source approximations
and target conditions, including the fixed nontrivial
signed value. As in LI.8, these are the chosen conditions
$\varphi(f)\in J$ and the required $f$ or frame membership
in the corresponding pronilpotent residuals; one does not
yet assert $D(k)=1$ or all earlier infinite constraints.
The finite target equations for the five $f$-factors are
among these prescribed conditions. They give the
finite-intersection property for the closed conditions
on the compact group of genuine gentle elements.
Compactness, with the original normalization conditions
retained in the same family, then produces one actual
$k$ with $\Delta(k)=1$, $D(k)=1$, every earlier retained
property, and all the new prime-power conditions
simultaneously.

The folded pentagon is one by the two-cycle relation, so
its defect belongs to $C$. At every odd prime the full
joint target kills the defect, and FF.6 gives
$\mathcal P(f_k)\in D_4'D_4^{\ell^n}$. The prime-two
containment follows from (FF.7d). Finally
\[
             \bigcap_{\ell\ {\rm prime},\ n\geq1}
                    D_4'D_4^{\ell^n}=D_4',
 \tag{FF.7e}
\]
because $D_4/D_4'\simeq\widehat{\mathbb Z}^{4097}$ is
separated by these finite prime-power quotients. This proves
(FF.7a). The same argument proves factorwise centralization
of $C/D_4'$: at odd primes a factor has identity joint
chain action, so every one of its commutators with $C$ is
killed by the faithful restriction in (FF.6a); at two use
the retained source action, and then use (FF.7e).

\paragraph{FF.8 (Proved): a genuine automorphism fixing the full fourth-power Ferrari kernel.}
For the same actual $k$ there is a continuous automorphism
\[
 \boxed{\quad \Phi_4\in\operatorname{Aut}(L/D_4'),
        \qquad \Phi_4|_{C/D_4'}=\mathrm{id},
        \qquad\Delta(k)=1.\quad}
 \tag{FF.8a}
\]
First the profinite kernel $C/D_4'$ is centreless. Let
$Q_4=C/D_4$, of order $8192$, and let $\Lambda$ be the
integral cycle lattice of its Cayley graph on the three
involutions. This graph has $8192$ vertices and $12288$
unoriented edges, so $\operatorname{rank}\Lambda=4097$.
For $q\ne1$ in $Q_4$, the vertex action has trace zero,
and a fixed unoriented edge is reversed. The rational
homology trace is therefore
\[
         \operatorname{tr}(q\mid\Lambda\otimes\mathbb Q)
                     =1-\#\{q\text{-fixed edges}\}\leq1,
\]
strictly below $4097$. This is the homology of an explicit
finite integral graph; no rational homology of a profinite
group is being used. Thus the integral deck action is
faithful. Its reduction modulo any odd prime is faithful
too: a finite two-power-order matrix congruent to the
identity modulo that prime belongs to a pro-odd-prime
congruence kernel and hence is the identity.

The signed edge module has no $Q_4$-invariants over
$\mathbb Q$, because the edge-stabilizing involution
reverses orientation. The rational cycle lattice also
has no invariants. Consequently the integral matrix
$\sum_{q\in Q_4}q$ is zero on $\Lambda$. It follows that
\[
                 (\widehat{\mathbb Z}^{4097})^{Q_4}=0:
\]
an invariant vector would satisfy $8192v=0$, and this
module is torsion-free. A central element of $C/D_4'$
first has trivial deck coordinate by the faithful action
on $\Lambda/3\Lambda$; its remaining coordinate is an
invariant in $D_4/D_4'$, and is therefore zero.

By FF.7 both braid conjugators
$g=f_k(a^2,b^2)$ and $h=f_k((ab)^3,c^2)$ centralize
$C/D_4'$. Apply the generator-and-relator argument of
FF.4 with the faithful inner action of this centreless
kernel in place of its finite joint representation.
The new generators have the old conjugation actions
on the kernel and the genuine gentle action on the
folded base. Every defining relator has trivial fold
and trivial kernel action, so is one. The resulting
homomorphism fixes the kernel pointwise, kills $D_4'$
before descending, and is an automorphism by the exact
sequence, since it is the identity on the kernel and
an automorphism on the full profinite base. This proves
(FF.8a). It is an actual relative source action and makes
no assertion that $k$ satisfies the full GT pentagon.
